\begin{table}[H]
\centering
\small
\caption{MFW i \amor{} jako komplementarne narzędzia analityczne}
\label{tab:mfw-vs-amor-fati}
\begin{tabularx}{\textwidth}{@{}L{0.20\textwidth}L{0.13\textwidth}>{\raggedright\arraybackslash}X@{}}
\toprule
Wymiar & Narzędzie & Ujęcie \\
\midrule
Pytanie badawcze & MFW & Jaki instrument redystrybucyjny poprawia równość i dobrobyt przy zadanym koszcie fiskalnym? \\
 & \amor{} & Jak BDP przechodzi przez bilanse, koszt kapitału, kurs, banki, inwestycje i decyzje firm o automatyzacji? \\
\addlinespace
Perspektywa czasu & MFW & Porównanie wariantów polityki w równowadze lub w analizie statycznej. \\
 & \amor{} & Ścieżka przejścia w horyzoncie 60 miesięcy. \\
\addlinespace
Jednostka BDP & MFW & Transfer na osobę, w polskim aneksie raportowany rocznie. \\
 & \amor{} & Miesięczny transfer na agenta reprezentującego gospodarstwo domowe, przeliczany następnie na ekwiwalent na osobę. \\
\addlinespace
Finansowanie & MFW & Warianty finansowania przez cięcia wydatków, VAT lub PIT; dla Polski koszt brutto w analizie statycznej. \\
 & \amor{} & Wariant deficytowy względem ścieżki bez BDP w \amor{}, z jawnym przejściem przez wydatki, deficyt i dług. \\
\addlinespace
Wyniki & MFW & Redystrybucja, ubóstwo, Gini, dobrobyt, koszt fiskalny. \\
 & \amor{} & Adopcja AI/hybrid, inwestycje, inflacja, stopa NBP, dług, deficyt, kurs, rachunek bieżący i bilanse banków. \\
\bottomrule
\end{tabularx}
\end{table}
